\begin{textsample}{NEG Dim 3 – 2016\_c6 – Score -5.00 – 2016\_c6/enade\_2016\_agro\_6f2e9fbda0.txt}  \label{ex:f3_neg_032}
CAPÍTULO 3 – ANÁLISE TÉCNICA DA PROVA Este capítulo tem por objetivo apresentar o desempenho dos estudantes concluintes de Agronomia no Enade/2016.
Para isso, foram calculadas as estatísticas básicas da prova em seu todo ( seção 3.1.1 ), bem como as estatísticas relacionadas aos Componentes de Formação Geral ( seção 3.1.2 ) e de Conhecimento Específico da Área ( seção 3.1.3 ).
Dadas as suas características, foram analisadas, em separado, as questões objetivas ( seção 3.2 ) e as questões discursivas ( seção 3.3 ).
Para as questões objetivas, foram disponibilizados os índices de facilidade e de discriminação ponto-bisserial, também em separado, para os Componentes de Formação Geral ( seção 3.2.1 ) e de Conhecimento Específico ( seção 3.2.2 ).
De cada componente, uma das questões foi escolhida para exemplificar a análise gráfica, relacionando as alternativas escolhidas pelos estudantes ( inclusive o gabarito ) com o número de acertos no componente.
O Anexo I apresenta a íntegra da análise gráfica para todas as questões objetivas.
Para cada uma das questões discursivas, os conteúdos dos tipos mais comuns de respostas dos estudantes são apresentados e comparados com o padrão de respostas esperado ( ver Anexo VIII com o padrão de respostas ).
Tomando como base as duas questões discursivas do Componente de Formação Geral, a seção 3.3.1.6 apresenta comentários sobre a correção das respostas com respeito à Língua Portuguesa.
Nas Tabelas 3.1 e 3.3 são apresentados o tamanho da população inscrita, ausente, presente e percentual de faltosos, respectivamente, por Grande Região e por Categoria Administrativa e Organização Acadêmica da IES.
As demais tabelas apresentam as seguintes estatísticas das notas15 : média do desempenho na prova, erro padrão da média, desvio padrão, nota mínima, mediana e nota máxima para cada um de seus componentes.
Tais estatísticas contemplam o total de estudantes concluintes da Área de Agronomia inscritos e presentes à prova do Enade/2016, tendo em vista agregações, ou por Grandes Regiões e o país como um todo, ou por Categoria Administrativa e Organização Acadêmica da IES.
Em relação aos gráficos de distribuição de notas, o intervalo considerado foi de 10 unidades, aberto à esquerda e fechado à direita, com exceção do primeiro intervalo, [ 0 ; 10 ], fechado em ambos os extremos.
Para os gráficos de distribuição das notas das 15 Uma definição dessas estatísticas pode ser encontrada no Glossário. 3.1 - CONDIÇÃO DE PRESENÇA E ESTATÍSTICAS BÁSICAS DA PROVA Esta seção apresenta a condição de presença, além das estatísticas selecionadas e histogramas da nota geral ( 3.1.1 ) e de cada componente : Formação Geral ( 3.1.2 ) e Conhecimento Específico ( 3.1.3 ).
São também apresentadas estatísticas selecionadas de subpopulações caracterizadas por Grande Região, Categoria Administrativa e Organização Acadêmica. 3.1.1 - Condição de Presença e Estatísticas Básicas Gerais A Tabela 3.1 apresenta a Condição de Presença na prova, por Grande Região, dos estudantes concluintes de Agronomia.
A população total de Inscritos foi de 12.175.
Destes, 11.198 estiveram Presentes, sendo 8,0 \% o índice de não comparecimento.
A Região de maior ausência foi a Centro-Oeste ( 11,3 \%, ) e a de menor ausência foi a Norte ( 6,2 \% ).
Já a Tabela 3.2 apresenta as Estatísticas Básicas da prova, por Grande Região, dos estudantes concluintes de Agronomia.
A Média das notas da prova como um todo ( nas seções seguintes serão analisados os Componentes de Formação Geral e de Conhecimento Específico ) foi 57,2, sendo que os alunos da região Norte obtiveram a Média mais baixa ( 52,7 ), e os da região Sul obtiveram a Média mais alta ( 61,0 ).
As demais médias foram : 54,4 na região Nordeste, 57,3 na região Sudeste e 55,1 na região Centro-Oeste.
O Desvio padrão para o Brasil como um todo foi 13,1, sendo o maior Desvio padrão encontrado na região Nordeste ( 13,5 ) e o menor, na região Sul ( 11,8 ), indicando uma dispersão um pouco menor das notas desta última região.
A região que obteve a maior nota Máxima foi a Sudeste ( 90,6 ), ao passo que a região que atingiu a menor nota Máxima foi a Norte ( 83,3 ).
A Mediana do Brasil como um todo foi 58,5, sendo a maior Mediana obtida na região Sul ( 62,2 ), e a menor obtida na Norte ( 53,5 ).
Em todas as regiões a nota Mínima foi zero.
Considerando -se as notas segundo Grande Região, observa -se que existe diferença estatisticamente significativa ao nível de 95 \% entre a menor Média, obtida na região Norte ( 52,7 ), e as médias de todas as demais regiões17.
O comportamento das notas dos estudantes de todo o Brasil pode ser observado no Gráfico 3.1 que apresenta um histograma com a distribuição das mesmas.
Essa é uma distribuição unimodal com a moda no intervalo ( 60 ; 70 ], seguido bem de perto pelo intervalo ( 50, 60 ].
A Tabela 3.3 apresenta informações referentes à Condição de Presença desagregadas por Categoria Administrativa e por Organização Acadêmica.
Da população total de Inscritos, 4.559 são de IES Privadas e 7.616 de IES Públicas.
Em relação à Organização Acadêmica, a maior participação foi obtida por estudantes de Universidades ( 8.572 ), seguido de Faculdades ( 2.322 ) e Centros Universitários ( 1.281 ).
A Categoria Administrativa de maior ausência foi a Privada ( 8,8 \% ), e entre as Organizações Acadêmicas foi a dos Centros Universitários ( 8,4 \% ), que ficou acima da taxa de ausência nacional de 8,0 \%.
Já a Tabela 3.4 apresenta informações referentes à Média da nota final desagregadas por Categoria Administrativa, e por Organização Acadêmica.
A Média das notas da prova como um todo foi 57,2.
Em relação à Categoria Administrativa, os estudantes das IES Públicas obtiveram Média mais alta ( 58,9 ), e os das IES Privadas obtiveram Média mais baixa ( 54,2 ), que a Média nacional.
Observa -se que existe diferença estatisticamente significativa entre as médias das notas das IES Públicas e Privadas.
A diferença entre as médias das regiões Sul e Norte ( 8,3 ), a maior e a menor Média, é superior à diferença entre IES Públicas e Privadas ( 4,7 ), caracterizando uma maior diversidade regional do que administrativa.
O Desvio padrão para as IES Privadas ( 13,2 ) foi superior ao do Brasil como um todo ( 13,1 ), indicando uma dispersão um pouco maior das notas nesta Categoria Administrativa.
No tocante à Organização Acadêmica, somente as Universidades obtiveram Média mais alta que a nacional ( 58,7 ).
A média dos Centros Universitários e das Faculdades foram 55,2 e 52,6, respectivamente.
Constata -se que existe diferença estatisticamente significativa ao nível de 95 \% nas médias das notas dos estudantes dos três tipos de Organização Acadêmica. 3.1.2 - Estatísticas Básicas no Componente de Formação Geral A Tabela 3.5 apresenta as Estatísticas Básicas em relação ao componente da prova que avalia a Formação Geral dos estudantes concluintes.
Os alunos de todo o Brasil obtiveram desempenho médio de 47,7.
Quanto à variabilidade, o Desvio padrão das notas dos estudantes do Brasil, como um todo, foi 15,8.
A maior Média foi obtida na região Sul ( 49,6 ), e a menor, na região Centro-Oeste ( 45,2 ).
As demais médias foram : 45,9 na região Norte, 48,1 na região Nordeste e 47,8 na região Sudeste.
Já o maior Desvio padrão foi obtido na região Nordeste ( 16,2 ), e o menor, na região Sul ( 15,2 ).
Os demais desvios padrões foram : 15,6 na região Norte, 15,9 na região Sudeste, e 15,8 da região Centro-Oeste.
A maior nota no Componente de Formação Geral da prova do Enade/2016 foi 91,8, obtida por, pelo menos, um aluno na região Sul.
A menor nota Máxima foi obtida na região Sudeste ( 87,4 ).
A Mediana do Brasil, como um todo, foi 48,5, sendo a menor Mediana encontrada na região Centro-Oeste ( 45,5 ), e a maior, encontrada na região Sul ( 49,9 ).
A nota Mínima nesta parte foi zero em todas as regiões, sem exceção.
Considerando -se as notas segundo Grande Região, observa -se que existe diferença estatisticamente significativa entre a maior Média das notas do Componente de Formação Geral, obtida na região Sul ( 49,6 ), e a Média das demais regiões.
O Gráfico 3.2 propicia a avaliação do desempenho dos estudantes do Componente de Formação Geral, a partir do histograma da distribuição das notas correspondentes.
A distribuição é unimodal, com moda em ( 40 ; 50 ], enquanto na prova, como um todo ( Gráfico 3.1 ), a moda foi alcançada no intervalo ( 60 ; 70 ].
Nota -se, ainda que, no Gráfico 3.2, as notas apresentam maior dispersão do que no Gráfico 3.1 ( distribuição das notas da prova ), confirmado pela comparação dos desvios padrões : 13,1 para a nota da prova como um todo e 15,8 para o Componente de Formação Geral.
Na Tabela 3.6 são apresentadas as informações referentes ao desempenho dos concluintes do Componente de Formação Geral, em diferentes agregações : Categoria Administrativa e Organização Acadêmica.
Observa -se que a maior Média foi obtida por estudantes de IES Públicas ( 50,2 ), com uma diferença estatisticamente significativa da obtida por estudantes de IES Privadas ( 43,6 ).
Nota -se que também há diferença estatisticamente significativa entre as médias dos três tipos de Organização Acadêmica, valor maior para Universidades ( 49,6 ) e menor para Faculdades ( 42,0 ). 3.1.3 - Estatísticas Básicas do Componente de Conhecimento Específico A Tabela 3.7 apresenta as Estatísticas Básicas referentes ao Componente de Conhecimento Específico da área de Agronomia.
A Média do desempenho dos alunos do Brasil, como um todo, foi 60,3.
A maior Média foi obtida na região Sul ( 64,8 ), e a menor, na região Norte ( 54,9 ).
As demais médias foram : 56,6 na região Nordeste, 60,5 na região Sudeste 58,4 na região Centro-Oeste.
Quanto à variabilidade das notas, o Desvio padrão do Brasil, como um todo, foi 14,6, sendo o maior Desvio padrão observado na região Nordeste ( 14,8 ), e o menor, na região Sul ( 13,1 ).
Os demais desvios foram : 14,6 nas regiões Norte e Sudeste e 14,7 na região Centro-Oeste.
A Mediana das notas dos estudantes de todo o Brasil foi 61,9.
A maior Mediana ocorreu na região Sul ( 66,2 ), e a menor, na região Norte ( 55,5 ).
As demais medianas foram : 58,1 na região Nordeste, 61,9 na região Sudeste e 59,8 na região Centro-Oeste.
A nota Máxima do Brasil, como um todo, foi 97,0, sendo obtida por, pelo menos, um aluno na região Sudeste.
As demais notas máximas foram : 88,7 na região Norte, 95,2 na região Nordeste, 94,2 na região Sul e 92,2 na região Centro-Oeste.
A nota Mínima foi zero em todas as regiões.
Observa -se que existe diferença estatisticamente significativa da menor Média das notas do Componente de Conhecimento Específico da região Norte ( 54,9 ) em relação às demais regiões, exceto em relação à região Nordeste.
Assim como os Gráficos 3.1 e 3.2, o Gráfico 3.3, apresentado a seguir, permite uma avaliação do desempenho de concluintes de Agronomia em relação ao Componente de Conhecimento Específico com um histograma da distribuição das notas correspondentes.
Esta também é uma distribuição unimodal, e o grupo modal é o ( 60 ; 70 ], o mesmo grupo modal da prova como um todo e acima do grupo modal para a Formação Geral.
A Tabela 3.8 apresenta uma comparação dos resultados em relação à Categoria Administrativa e à Organização Acadêmica, agora levando em conta o desempenho de alunos do Componente de Conhecimento Específico da prova.
No que se refere à Organização Acadêmica, a maior Média foi das Universidades ( 61,7 ), vindo a seguir a dos Centros Universitários ( 58,4 ) e, depois, a das Faculdades ( 56,2 ).
O menor Desvio padrão, e abaixo do valor para o Brasil, como um todo, foi o das Universidades ( 14,2 ).
As Universidades também obtiveram as maiores notas Máxima ( 97,0 ) e Mediana ( 63,0 ).
Os Centros Universitários obtiveram nota Máxima igual a 92,9 e as Faculdades obtiveram nota Máxima igual a 92,8.
As medianas foram : 59,3 nos Centros Universitários e 57,1 nas Faculdades.
A nota Mínima foi zero para as Universidades e Centros Universitários, e 7,1 para as Faculdades Observa -se que existem diferenças estatisticamente significativas ao nível de 95 \% no Componente de Conhecimento Específico entre as notas das Organizações Acadêmicas.
Quanto à Categoria Administrativa, observa -se um comportamento semelhante àquele da parte de Formação Geral e à prova como um todo, ou seja, existe diferença estatisticamente significativa entre as médias das IES Públicas ( 61,8 ) e IES Privadas ( 57,7 ).
Neste caso também, a maior Média foi obtida por alunos de IES Públicas de ensino. 3.2 - ANÁLISE DAS QUESTÕES OBJETIVAS Esta seção apresenta estatísticas selecionadas e histogramas das Questões Objetivas de Formação Geral ( 3.2.1 ).
São também apresentadas e comparadas as médias das subpopulações caracterizadas por Grande Região. 3.2.1 - Componente de Formação Geral A Tabela 3.9 apresenta as Estatísticas Básicas relativas às oito questões objetivas do componente da prova que abrange a Formação Geral dos estudantes.
A Média do Brasil foi 42,0.
A menor Média foi encontrada na região Centro-Oeste ( 38,7 ), e a maior, na região Sul ( 43,9 ).
As demais médias foram : 40,2 na região Norte, 42,5 na região Nordeste e 42,3 na região Sudeste.
O Desvio padrão do Brasil foi 20,0, sendo o maior Desvio padrão encontrado nas regiões Nordeste e Sudeste ( 20,1, cada ), e o menor, na região Norte ( 19,4 ).
Os demais desvios foram : 19,9 na região Sul e 19,8 na região Centro-Oeste.
As medianas do Brasil como um todo e de todas as regiões foi 37,5.
As notas máximas ( 100,0 ) e as notas mínimas ( 0,0 ) foram iguais para todas as regiões.
A Tabela 3.10 apresenta o Índice de Facilidade e o Índice de Discriminação ( ponto-bisserial ) para cada uma das questões objetivas do Componente de Formação Geral.
Quanto ao Índice de Facilidade, foram usadas as seguintes cores para diferenciar o nível de dificuldade da questão : • Azul para as questões classificadas com índice Muito fácil ( \&gt ; =0,86 ), verde para as questões classificadas com índice Fácil ( 0,61 a 0,85 ), amarelo para as questões classificadas com Médio ( 0,41 a 0,60 ), vermelho para as questões classificadas com Difícil ( 0,16 a 0,40 ) e roxo para as questões classificadas com Muito difícil ( \&lt ; =0,15 ).
Já quanto ao Índice de Discriminação, foram usadas as seguintes cores para 0,29 ), as classificadas com Bom receberam a cor verde ( 0,30 a 0,39 ) e as classificadas com Muito bom ( \&gt ; =0,40 ) receberam a cor azul.
As questões objetivas do Componente de Formação Geral, segundo o Índice de Facilidade, foram assim avaliadas : das oito questões, nenhuma teve o Índice de Facilidade classificado como Muito fácil e uma questão foi tida como Fácil, por ter índice de acertos 0,78.
Quatro questões foram consideradas com índice de dificuldade Médio, situando -se no intervalo entre 0,41 e 0,60 do Índice de Facilidade, ou seja, houve entre 41,0 \% e 60,0 \% de acertos, enquanto duas questões foram classificadas na categoria Difícil, situando -se no intervalo entre 0,16 e 0,40.
Por fim, uma questão apresentou 15 \% de acertos, razão pela qual seria classificada como Muito difícil.
Como já comentado, para análise das questões objetivas, relativas à Formação Geral, segundo o poder de discriminação, utilizou -se o Índice de Discriminação ( ponto-bisserial ).
Nesta análise, as questões foram assim avaliadas : seis das oito questões apresentaram índice acima ou igual a 0,40 e, assim, foram classificadas com índice Muito bom para esse grupo de estudantes e duas questão tiveram Índice de Discriminação Bom, com valor entre 0,30 e 0,39.
Nenhuma questão teve nível Fraco ou Médio de discriminação para esse grupo de estudantes.
O Índice de Facilidade variou de 0,15 a 0,78, e o de Discriminação, de 0,32 a 0,51.
As seis questões com Índice de Discriminação Muito bom figuraram entre níveis de dificuldade Fácil, Médio e Difícil : uma classificada na categoria Fácil ( a questões 6 ) do Índice de Facilidade, quatro, na categoria Médio ( questões 1, 2, 3 e 7 ) e uma na categoria Difícil ( questão 4 ).
Em particular, a questão 2 foi a que apresentou o maior poder discriminatório, com índice 0,53, e foi considerada Médio em termos de facilidade, com uma proporção de 0,45 acertos.
O máximo de acertos foi alcançado pela questão 6 com um Índice de Facilidade de 0,78.
A questão de número 5 apresentou um Índice de Facilidade de 0,15, ou seja, menos do que um quinto dos estudantes conseguiu resolvê -la, dentro do universo de participantes.
Seu Índice de Discriminação foi Bom ( 0,32 ).
O Gráfico 3.4, para exemplificar, analisa o comportamento da questão de número 2 de Formação Geral.
Trata -se de uma questão considerada Médio, com relação à facilidade e a que obteve o maior Índice de Discriminação dessa parte da prova.
Neste gráfico, cada uma das cinco curvas representa o percentual de respostas em determinada alternativa da questão, em função do número de acertos dos estudantes nessa parte da prova ( Formação Geral/Múltipla Escolha ), antes de possíveis eliminações pelo critério do ponto-bisserial.
A curva em azul corresponde à alternativa A, a alternativa correta para esta questão.
Observa -se que entre os estudantes com menor número de acertos, nessa parte do exame, a situação mais frequente foi a escolha da alternativa E ( em vermelho ), incorreta.
Por exemplo, entre os estudantes que acertaram três questões, 33,0 \% escolheram a alternativa B ( verde ), 26,7 \% escolheram a alternativa E ( em vermelho ), 1,8 \% escolheram a alternativa C ( preto ) e 1,1 \%, a D ( roxo ).
Entre os que acertaram três respostas entre as questões de múltipla escolha de formação geral, 0,3 \% marcou mais de uma alternativa, invalidando a questão, e nenhum aluno deixou a questão em branco.
À medida que o número de acertos aumenta, indicando desempenho melhor nessa parte da prova, aumenta concomitantemente a proporção de estudantes que selecionaram a alternativa correta A, atingindo 100 \% para os estudantes com 8 acertos.
Essa análise permite verificar como a questão discriminou os grupos de desempenho, justificando o alto índice obtido na questão.
Cumpre notar que não é possível inferir deste gráfico nem o índice de facilidade ( que seria uma média da proporção ponderada pela quantidade de alunos com cada uma das notas ), nem o índice de discriminação ponto-bisserial, por razão equivalente.
No caso extremo no qual a grande concentração dos acertos dos alunos fosse abaixo de 4, o índice de facilidade seria obrigatoriamente abaixo de 40 \% ( neste exemplo ).
Caso a concentração fosse em 6 acertos ou mais, o índice seria obrigatoriamente acima de 60 \%.
Os gráficos relativos às demais questões de Formação Geral constam do Anexo I. 3.2.2 - Componente de Conhecimento Específico A Tabela 3.11 apresenta as Estatísticas Básicas em relação às questões objetivas do Componente de Conhecimento Específico da prova de Agronomia por Grande Região.
A Média do Brasil deste componente foi de 61,8.
A menor Média foi observada na região Norte ( 56,1 ), e a maior, na região Sul ( 66,6 ).
O Desvio padrão de todo o Brasil foi 15,7, sendo o menor Desvio padrão encontrado na região Sul ( 14,2 ), e o maior, nas regiões Nordeste e Centro-Oeste ( 15,8, cada ).
A Mediana de todo o Brasil foi 62,5, a mesma encontrada em duas das regiões : Sudeste e Centro-Oeste.
Na região Norte a Mediana foi menor ( 54,2 ), na região Nordeste foi 58,3 e na região Sul, maior ( 66,7 ).
A nota Máxima da prova foi 100,0, obtida nas questões objetivas do Componente de Conhecimento Específico, por, pelo menos, um aluno nas regiões Sudeste e Sul.
Nas demais regiões, a nota Máxima da prova foi : 91,7 na região Norte e 95,8 nas regiões Nordeste e Centro-Oeste.
Em todas as regiões a nota Mínima foi zero.
A Tabela 3.12 apresenta os Índices de Facilidade e Discriminação ( ponto-bisserial ) das questões objetivas do Componente de Conhecimento Específico para os estudantes de Agronomia.
Para facilitar a diferenciação das questões, usaremos as mesmas cores da Tabela 3.10 para as diferentes classificações dos Índices de Facilidade e de Discriminação.
A partir do Índice de Facilidade obtido, pode -se concluir que somente duas questões objetivas da prova foram consideradas, pelo menos, Difícil.
Das 25 questões válidas ( as questões 17 e 21 foram anuladas pela banca ), mais da metade ( 15 ) foi classificada como Fácil, e outras oito consideradas como Médio.
Já quanto aos índices de discriminação das questões objetivas do Componente de Conhecimento Específico da prova, tem -se como resultado a seguinte classificação : 19 das 25 questões válidas foram consideradas como boas, enquanto uma delas tiveram Índice de Discriminação Muito bom.
Assim, para exatamente quatro quintos das questões – 20 em 25 – os Índices de Discriminação foram Bom ou Muito bom.
Dentre as demais, quatro delas foram classificadas como Médio e uma como Fraco, sendo cinco, por conseguinte, a quantidade de questões nos dois patamares mais baixos de discriminação.
Constata -se, assim, que a prova – no que se refere ao Componente de Conhecimento Específico – possuía capacidade média de discriminar entre aqueles que dominam ou não o conteúdo.
Dentre as questões que alcançaram os maiores índices de discriminação, uma delas, a de número 29, foi classificada com Índice de Discriminação Muito bom, situando -se no intervalo de 0,40 a 0,50 do índice, e esta foi classificada como Fácil, quanto ao Índice de Facilidade.
A questão de número 12 foi a mais difícil dentre as 25 questões válidas específicas, com baixo Índice de Facilidade, com 0,24 \% de acertos.
Essa questão apresentou poder discriminatório igualmente baixo, 0,11, o que comprova ter sido esta questão a mais difícil para os estudantes.
Tal questão foi, portanto, pelo critério ponto-bisserial, considerada inadequada.
Por isso, a questão 12, além das questões 17 e 21, foi eliminada do cômputo da nota final.
A título de exemplo das análises do comportamento das questões objetivas, o Gráfico 3.5 analisa a questão 29 do Componente de Conhecimento Específico.
Esta questão foi considerada pelas respostas dos estudantes avaliados na prova como uma questão com nível de facilidade Fácil, 0,68, ou seja, 68,0 \% dos estudantes assinalaram acertadamente a opção E, correspondente ao gabarito.
Seu Índice de Discriminação foi igual a 0,46, classificado como Muito bom, também sendo esta questão a que apresentou o maior valor.
Neste gráfico, cada uma das cinco curvas representa o percentual de respostas em determinada alternativa da questão 29, em função do número de acertos dos estudantes nessa parte da prova, antes de possíveis eliminações de questões pelo critério do ponto-bisserial.
A alternativa correta E, representada no gráfico pela curva em vermelho, foi escolhida em maiores proporções pelos alunos com desempenho melhor nessa parte da prova.
Já as alternativas incorretas, também denominadas distratores, foram selecionadas, principalmente, por aqueles com notas mais baixas.
Neste caso, também a soma não é sempre 100 \% por causa das questões não respondidas ou com mais de uma opção marcada.
Aqueles com nota zero, na sua quase totalidade deixaram esta questão em branco ou marcaram mais de uma alternativa, comportamento considerado inválido.
A proporção de alunos que selecionou a resposta correta E aumenta gradativamente, chegando a atingir 100 \% para 24 acertos ( com uma pequena queda para os que acertaram 23 ), enquanto a proporção dos que escolheram alternativas incorretas decai, a partir de sete acertos, como função do número de acertos nessa parte da prova, com exceção da alternativa B ( em verde ), que apresenta um crescimento no último intervalo.
Os gráficos relativos às demais questões do Conhecimento Específico constam do Anexo I. 3.3 - ANÁLISE DAS QUESTÕES DISCURSIVAS Esta seção apresenta estatísticas selecionadas e histogramas das Questões Discursivas de Formação Geral ( 3.3.1 ) e Conhecimento Específico ( 3.3.2 ).
São também apresentadas e comparadas as médias de subpopulações caracterizadas por Grande Região. 3.3.1 - Componente de Formação Geral As análises dos resultados de desempenho dos estudantes de Agronomia nas duas questões discursivas relativas à Formação Geral encontram -se na Tabela 3.13 e no Gráfico 3.6.
Na Tabela 3.13, observa -se que a nota Média nesse conjunto de questões foi acima da obtida nas objetivas.
Os estudantes, de todo o Brasil, obtiveram, em Formação Geral, Média 42,0 nas questões objetivas e 56,4 nas questões discursivas.
Pode -se notar também que o Desvio padrão nesse conjunto de questões foi um pouco menor do obtido nas objetivas, 20,0 nas questões objetivas e 19,7 nas questões discursivas.
A maior Média foi obtida na região Sul ( 58,1 ), e a menor, na região Norte com 54,5.
A Mediana de todo o Brasil, neste componente, foi 60,0.
Nas regiões Norte e Centro-Oeste a Mediana foi menor, 59,0 em ambas.
Nas demais regiões foi igual ou maior que o valor para o Brasil : Nordeste e Sul ( 61,0 em ambas ) e Sudeste ( 60,0 ).
A nota Máxima ( 96,0 ) foi a mesma nas regiões Nordeste, Sudeste e Sul, sendo 94,0 na região Norte e 95,0 na região Centro-Oeste.
A nota Mínima ( 0,0 ) foi a mesma em todas as regiões do Brasil.
O Gráfico 3.6 representa a distribuição das notas nas questões discursivas do Componente de Formação Geral.
A moda desta distribuição ocorre no intervalo ( 60 ; 70 ].
Destaca -se também o intervalo [ 0 ; 10 ] com distribuição próxima a 6 \% do total de notas, sendo que no intervalo [ 0 ; 10 ] incluem -se além da nota zero, a frequência de alunos que deixaram este tipo de questão em branco.
Na sequência, os resultados verificados para cada uma das questões discursivas de Formação Geral serão apresentados, estabelecendo -se relações com os conteúdos abordados em cada uma delas.
Os comentários da Banca de docentes corretores, a respeito do observado na correção das respostas dos estudantes, suas impressões e conclusões serão apresentados junto à análise de cada questão.
Cumpre esclarecer que, tendo em vista que as questões discursivas de Formação Geral são padronizadas, ou seja, constam de todas as provas, os comentários da Banca são os mesmos para todas as carreiras acadêmicas, sendo direcionados a todos os estudantes que participaram do Enade/2016.
A seguir, serão analisados os desempenhos da Área de Agronomia nas duas questões discursivas de Formação Geral do Enade/2016, comparando -se os resultados obtidos com comentários para cada questão. 3.3.1.1 - Análise de Conteúdo da Questão Discursiva 1 do Componente de Formação Geral Os dados de Agronomia, obtidos a partir das respostas à questão 1, encontram -se na Tabela 3.14 e no Gráfico 3.7.
Nessa questão – de melhor desempenho dentre as duas de Formação Geral – os alunos, de todo o Brasil, tiveram Média 66,1.
A maior Média para a questão 1 foi obtida na região Sul ( 68,3 ), e a menor, na região Norte ( 63,9 ).
Quanto à variabilidade das notas, o Desvio padrão de todo o Brasil foi 27,5.
O menor Desvio padrão foi obtido na região Sul ( 25,3 ), e o maior Desvio padrão foi obtido na região Norte ( 29,8 ).
As medianas de todas as regiões e do Brasil como um todo foram iguais ( 70,0 ).
As notas máximas e mínimas da questão discursiva 1 foram, sem exceção, as mesmas para todas as regiões do Brasil, respectivamente, 100,0 e 0,0.
O Gráfico 3.7 mostra a distribuição das notas na questão discursiva 1 do Componente de Formação Geral.
Observa -se que a maior frequência corresponde aos alunos que obtiveram nota dentro do intervalo ( 70 ; 80 ], correspondendo à moda da distribuição.
Destacam -se, também, os alunos que deixaram a questão em branco com aproximadamente 7,0 \% do total. 3.3.1.2 - Comentários sobre a correção de Conteúdo das respostas à Questão Discursiva 1 A questão discursiva 1 alinha -se aos conteúdos programáticos previstos nas diretrizes curriculares, quanto às habilidades e às competências para a atualização permanente e aos conhecimentos sobre a realidade brasileira, mundial, tangenciando outras áreas do conhecimento.
O enunciado é claro e indica um tema de extrema relevância, qual seja, as migrações humanas e todo o complexo de problemas que esse movimento causa, quer na origem, quer no seu ponto de destino.
A questão é extremamente atual, maximizada pelos acompanhamentos midiáticos que espetacularizam cenas dramáticas do processo de fuga de massas humanas das guerras, como a morte de imigrantes.
Além do aspecto internacional, deve se ter em conta que o Brasil sofreu vários influxos migratórios, desde o forçado, pela transferência de escravos da África, até o incentivado, dos italianos, alemães e outros europeus para trabalhar na agricultura.
Isso sem considerar a própria colonização portuguesa originária e as migrações posteriores dos mesmos portugueses.
Posteriormente, o Brasil tornou -se um país de emigrantes que foram \textbf{aportar} a outros locais na busca de melhor situação econômica.
Muitos migraram para a América do Norte e para a Europa, principalmente.
Mais recentemente, por força do crescimento econômico pátrio, novas levas de imigrantes desembarcaram no Brasil, muitos deles provenientes do Haiti, embora tenhamos recebidos vizinhos latinos, como bolivianos, paraguaios, chilenos, argentinos, venezuelanos, peruanos, colombianos que também vieram em busca de uma vida melhor.
Como se verifica, a temática indicada no enunciado pertence à realidade conhecida amplamente pelos estudantes, pois, também divulgada por todos os meios de comunicação.
Trata -se de tema geral que, normalmente, não se inclui nas grades curriculares superiores, mas é de conhecimento geral e relacionado aos direitos humanos, este sim, ponto de integração às diretrizes curriculares.
A linguagem utilizada foi clara, indicando a existência dos fatos vinculados aos movimentos migratórios.
Por outro lado, o texto demanda uma reflexão sobre a realidade dos indivíduos que são compelidos ao afastamento dos seus lares por motivos externos, como guerras, conflitos civis, religiosos, perseguições das mais variadas, questões econômicas.
Como se trata de texto contendo uma ampla gama de possibilidades, é natural que ocorram interpretações diversas das esperadas.
Do mesmo modo, a contemporaneidade do tema encaminha o enunciado para ser qualificado como questão de média dificuldade, em comparação com as demais questões apresentadas.
O padrão de resposta indicou possíveis interpretações do texto.
No que concerne ao grau de profundidade, mostrou a necessidade de conhecimento de vários itens pertinentes à migração e que são refletidos nas várias possibilidades de resposta previstas pelo padrão.
Pela análise das provas realizadas, constatou -se que muitos foram os que identificaram os itens indicados no padrão de resposta.
Essa circunstância poderia ser justificada, tendo em vista que o tema vem sendo bastante difundido na mídia social, gerando elementos naturais para a resposta, bem como da história do país, que, intimamente, vincula -se ao processo migratório, quer ativo, remetendo pessoas para outros países, quer passivo, na recepção de pessoas.
Essa característica também permite responder aos termos postos quanto aos Estados que recebem os migrantes.
Assim, a indicação dos aspectos positivos e negativos também é bastante difundida.
As respostas medianas e boas apreenderam o sentido do enunciado, e indicaram os motivos expressos no padrão de respostas, como os ensejadores das imigrações, notadamente, para a Europa, que seria o foco do texto.
Por outro lado, também identificaram os problemas dos imigrantes, no mesmo padrão indicado para as respostas desejadas.
Alguns concluintes procuraram exemplos de comparação na realidade brasileira, como os clássicos imigrantes italianos, alemães e de outras nacionalidades, bem como os atuais haitianos e demais povos da América Latina.
Uma minoria não apreendeu o sentido do enunciado e, somente, apresentou elementos da realidade nacional.
Esses foram avaliados no conjunto de provas fracas.
As respostas demonstraram que o concluinte está atualizado, em regra, com o tema da imigração, indicando os fundamentos dessa movimentação populacional que, nos últimos tempos, tem se agravado pela incidência de conflitos armados em vários países, ensejadores de reflexos na segurança dos habitantes desses locais, bem como o surgimento de focos de doenças e ausência dos serviços públicos básicos, além da fome, decorrente da impossibilidade de produzir alimentos ou mesmo comerciá -los.
Do ponto de vista dos Estados que recebem o fluxo de imigrantes, também restou clara a percepção das dificuldades vinculadas à ausência de estrutura quanto aos serviços públicos ( saúde, educação, segurança, moradia, alimentação ) para receber um número expressivo de imigrantes.
Em época de recessão econômica, a presença de estrangeiros acarreta a competição pela busca de empregos, gera movimentos xenofóbicos, bem como choques culturais e políticos.
Deve -se ter em mira que esses movimentos têm um viés de insegurança, pelos países receptores, diante da crise causada por movimentos considerados terroristas, embora a imensa maioria não tenha qualquer vínculo com esses movimentos.
Os recentes atentados na Europa indicam que os seus autores têm origem em comunidades de imigrantes, embora muitos já tenham nascido nos países da comunidade europeia.
Essas características negativas, no entanto, paradoxalmente, convivem com aspectos positivos, o que encaminha a questão para a necessidade dos jovens imigrantes, que ocuparão postos de trabalho não desejados pelos habitantes locais, bem como, compensarão os baixos índices de crescimento populacional.
Existe outro reflexo positivo que está relacionado ao sistema previdenciário, que terá novos contribuintes nos trabalhadores estrangeiros que se integram ao sistema.
Esses pontos foram abordados e indicam que existe uma consciência dos concluintes dos cursos universitários quanto aos principais itens relacionados ao tema que foi posto para análise. 3.3.1.3 - Análise de Conteúdo da Questão Discursiva 2 do Componente de Formação Geral A Tabela 3.15 mostra que o desempenho médio dos estudantes na questão discursiva 2 ( média 44,8 ) foi inferior ao obtido na questão discursiva 1 ( média 66,1 ).
A região Sul foi aquela cuja Média, também nessa questão, foi maior ( 46,0 ), e a de menor Média foi a região Norte ( 43,7 ).
Quanto à variabilidade das notas, o Desvio padrão de todo o Brasil foi 23,5, inferior ao obtido na questão discursiva 1 ( 27,5 ).
O maior desvio nessa questão foi obtido na região Nordeste ( 24,5 ), enquanto o menor foi obtido na região Sul ( 22,6 ).
A Mediana de todo o Brasil foi 50,0, a mesma de todas as regiões.
As notas máximas ( 100,0 ) e as notas mínimas ( 0,0 ) também foram as mesmas em todas as regiões, sem exceção.
O Gráfico 3.8 mostra a distribuição das notas na questão discursiva 2 do Componente de Formação Geral.
Observa -se que a maior frequência corresponde aos alunos que obtiveram nota dentro do intervalo ( 50 ; 60 ].
Destacam -se, também, os alunos que deixaram a questão em branco e os alunos que obtiveram nota dentro do intervalor ( 30 ; 40 ], caracterizando -se como máximos locais.
Nota -se, ainda, que as notas ficaram mais concentradas em comparação à questão discursiva de número 1.
Isto pode ser constatado também pela comparação do desvio padrão das notas da questão discursiva 2 ( 23,5 ) e da questão discursiva 1 ( 27,5 ). 3.3.1.4 - Comentários sobre a correção de Conteúdo das respostas à Questão Discursiva 2 A questão discursiva 2 tem seu conteúdo integrado às Diretrizes Curriculares, uma vez que a proteção aos direitos humanos e às relações de gênero constituem matrizes fundamentais para conhecimento dos profissionais de nível superior.
O enunciado é claro e indica um tema de extrema relevância, qual seja, a análise da questão da mulher submetida a injustas agressões e os seus reflexos pessoais, sociais e familiares.
Também enfrenta a busca pela maior proteção no concernente às agressões, com a instituição de um modelo legal mais repressivo e que foi gerado por uma situação especifica de ataque à mulher, com consequências dramáticas para a vítima, Maria da Penha, cuja lei levou seu nome. notas ficaram mais concentradas em comparação à questão discursiva de número 1.
Isto pode ser constatado também pela comparação do desvio padrão das notas da questão discursiva 2 ( 23,5 ) e da questão discursiva 1 ( 27,5 ). 3.3.1.4 - Comentários sobre a correção de Conteúdo das respostas à Questão Discursiva 2 A questão discursiva 2 tem seu conteúdo integrado às Diretrizes Curriculares, uma vez que a proteção aos direitos humanos e às relações de gênero constituem matrizes fundamentais para conhecimento dos profissionais de nível superior.
O enunciado é claro e indica um tema de extrema relevância, qual seja, a análise da questão da mulher submetida a injustas agressões e os seus reflexos pessoais, sociais e familiares.
Também enfrenta a busca pela maior proteção no concernente às agressões, com a instituição de um modelo legal mais repressivo e que foi gerado por uma situação especifica de ataque à mulher, com consequências dramáticas para a vítima, Maria da Penha, cuja lei levou seu nome.
Deve ser recordado que os conflitos vinculados às relações familiares sempre foram circunscritos ao seu âmbito e, quando dele extravasavam, por força de agressões contundentes à integridade física de um dos partícipes da relação, em geral a mulher, raramente redundavam em punições relevantes, sendo comum o pagamento em cestas básicas de alimentos, como sanção.
Após a edição da lei Maria da Penha ( Lei no. 11.340, sancionada em 7 de agosto de 2006 ), além das sanções sofrerem modificações, houve a criação de modelos de proteção e atendimento especial para os casos de violência doméstica, procurando proteger as mulheres agredidas antes de danos maiores, muitas vezes irreversíveis, incluindo aí o feminicídio, também incluído no Código Penal, como circunstância qualificadora do crime de \textbf{homicídio}, e incluí -lo no rol dos crimes hediondos.
A lei Maria da Penha busca realizar, como indicado no seu pórtico, a Convenção sobre a Eliminação de Todas as Formas de Discriminação contra as Mulheres e a Convenção Interamericana para Prevenir, Punir e Erradicar a Violência contra a Mulher.
No seu texto, várias são as formas de violência doméstica e familiar contra a mulher, identificadas como sendo, entre outras : “ a violência física, entendida como qualquer conduta que ofenda sua integridade ou saúde corporal ; a violência psicológica, entendida como qualquer conduta que lhe cause dano emocional e diminuição da \textbf{autoestima} ou que lhe prejudique e perturbe o pleno desenvolvimento ou que vise a degradar ou controlar suas ações, comportamentos, crenças e decisões, mediante ameaça, constrangimento, \textbf{humilhação}, manipulação, isolamento, vigilância constante, perseguição contumaz, insulto, chantagem, ridicularização, exploração e limitação do direito de ir e vir ou qualquer outro meio que lhe cause prejuízo à saúde psicológica e à autodeterminação ; a violência sexual, entendida como qualquer conduta que a constranja a presenciar, a manter ou a participar de relação sexual não desejada, mediante intimidação, ameaça, coação ou uso da força ; que a induza a comercializar ou a utilizar, de qualquer modo, a sua sexualidade, que a impeça de usar qualquer método contraceptivo ou que a force ao matrimônio, à gravidez, ao aborto ou à prostituição, mediante coação, chantagem, suborno ou manipulação ; ou que limite ou anule o exercício de seus direitos sexuais e reprodutivos ; a violência patrimonial, entendida como qualquer conduta que configure retenção, subtração, destruição parcial ou total de seus objetos, instrumentos de trabalho, documentos pessoais, bens, valores e direitos ou recursos econômicos, incluindo os destinados a satisfazer suas necessidades ; a violência moral, entendida como qualquer conduta que configure calúnia, difamação ou injúria. ” ( Conforme texto extraído do artigo sétimo da referida lei ).
Também se buscam medidas de integração entre os vários órgãos públicos, abrangendo repressão, proteção, conscientização, bem como utilização de equipes multidisciplinares, tendo em vista os graves danos, de variados matizes, que redundam das práticas agressivas.
Nessa linha, também existe na lei a previsão de locais de atendimento específico ( centros de atendimento integral e multidisciplinar para mulheres e respectivos dependentes em situação de violência doméstica e familiar ; casas-abrigo para mulheres e respectivos dependentes menores em situação de violência doméstica e familiar ; delegacias, núcleos de defensoria pública, serviços de saúde e centros de perícia médico-legal especializados no atendimento à mulher em situação de violência doméstica e familiar ; programas e campanhas de enfrentamento da violência doméstica e familiar ; centros de educação e de reabilitação para os agressores ).
O Poder Judiciário, também, deve especializar o seu atendimento para aplicar as medidas de restrição ao agressor de forma mais eficiente.
A linguagem utilizada na questão foi clara, indicando a grave violação dos direitos das mulheres injustamente agredidas e a luta pela busca de sanções adequadas aos agressores e de proteção pessoal e para os de sua família.
Ademais, o tema não está circunscrito ao Brasil, mas o país possui um número muito expressivo de atos que compõem esse cenário.
Por outro lado, o texto demanda uma reflexão sobre a realidade, e abarca diversos segmentos da sociedade, pois as situações de agressão não se restringem às camadas mais populares.
Como se trata de texto contendo uma ampla gama de possibilidades, é natural que ocorram interpretações diversas das esperadas.
O padrão de resposta identifica a projeção do elaborador na resposta a ser apresentada pelo estudante.
Está em linguagem correta e indica possíveis interpretações do texto.
Suas opções foram bastante amplas, abrangendo a maior parte das respostas encontradas.
No que concerne ao grau de profundidade, o padrão de resposta indicou a necessidade de conhecimento de lei especial que regula a violência doméstica, o que exige ou qualificação da área de Direito ou que os currículos dos cursos dos participantes desse ciclo do Enade tenham oferecido alguma disciplina introdutória da área jurídica.
Ao analisar as provas realizadas, constatamos que, na primeira parte, houve maior facilidade em responder aos itens indicados no padrão de resposta, pois atinentes, não só à formação profissional, mas, também, ao conhecimento da realidade, pois, infelizmente, a violência doméstica está disseminada na sociedade e nos vários pontos do país.
Mas, verificou -se que as respostas vieram, no mais das vezes, desorganizadas, dificultando a correção.
Apesar disso, restou claro que a maioria esmagadora das respostas confluiu para as indicações do padrão apresentado, sendo relevante indicar que o tema remeteu para a luta em relação à igualdade de gênero, temática candente em especial para as concluintes do sexo feminino.
Diversos textos se apresentaram de forma concatenada, fundamentada e, muitas vezes, lamentando que, nos dias atuais, ainda se discuta a necessidade de proteger as mulheres dessa espécie de violência.
Em muitos casos, também, os concluintes fizeram um liame entre o acesso à educação e a melhora das relações entre homens e mulheres, ao permitir que a conclusão de maior grau de educação formal conduzisse à diminuição do índice de violência contra a mulher.
Foram identificadas respostas consoantes com o padrão apresentado, indicando os itens que nele constam em relação aos reflexos pessoais, na família e na sociedade.
As respostas dos estudantes raramente faziam referência à organização protetiva instituída pela Lei Maria da Penha, como esperado pelo padrão de resposta.
A respeito de tal aspecto o padrão indicava como uma das possibilidades de respostas ( item ‘ b ’ ) : a criação de casas de refúgio ou casa-abrigo, para acolher mulheres vítimas de violência doméstica ; reforço às Delegacias de Atendimento à Mulher ; inclusão da vítima em programas assistenciais do governo, programas de proteção à vítima e à testemunha, transferência de local de trabalho, se a vítima for servidora pública.
Como paradigma, indicaram -se as Delegacias da Mulher que são organizadas, com a quase integralidade dos seus membros de policiais do sexo feminino, e especializadas no atendimento a mulheres submetidas a situações de violência.
Essa preocupação gerencial tem por escopo bloquear os preconceitos que acompanham essas dramáticas histórias de violência, o que gera uma incompreensão e uma ineficiência na rede de proteção.
As respostas demonstram que os concluintes têm ciência do problema, da sua gravidade, do vínculo com a questão de gênero e com a submissão da mulher como efeito da sociedade patriarcal, ainda que tratado de forma superficial, com base na mídia e, algumas vezes, com apoio de experiências pessoais.
Identificou -se uma generalizada compreensão de que o problema está sendo enfrentado e que ainda existe um longo caminho no percurso para a sua resolução. 3.3.1.5 - Análise de Língua Portuguesa das Questões Discursivas do Componente de Formação Geral Os dados de Agronomia, obtidos a partir das respostas às questões discursivas do Componente de Formação Geral, no que tange à Língua Portuguesa, encontram -se na Tabela 3.16 e no Gráfico 3.9.
Nesse aspecto, os alunos, de todo o Brasil, tiveram Média 60,1.
A maior Média com respeito à Língua Portuguesa foi obtida na região Sul ( 61,9 ), e a menor, na região Norte ( 57,7 ).
Quanto à variabilidade das notas, o Desvio padrão de todo o Brasil foi 18,3.
O menor Desvio padrão foi obtido na região Sul ( 16,2 ) e o maior Desvio padrão foi obtido na região Norte ( 19,8 ).
A Mediana das notas de Língua Portuguesa foi 65,0 para as regiões Nordeste, Sudeste e Sul, e 62,5 para as regiões Norte e Centro-Oeste.
A nota Máxima para todo o Brasil foi de 95,0, com, pelo menos, um aluno tirando essa nota na região Nordeste.
As demais regiões obtiveram nota Máxima 90,0 ( regiões Norte, Sudeste e Centro-Oeste ) e 92,5 ( região Sul ).
Além disso, a nota Mínima foi zero em todas as regiões do país.
O Gráfico 3.9 mostra a distribuição das notas de Língua Portuguesa do Componente de Formação Geral.
Observa -se que a maior frequência ( mais de 35,0 \% ) corresponde aos alunos que obtiveram nota no intervalo ( 60 ; 70 ].
Destaca -se, também, os alunos que deixaram ambas as questões em branco representando quase 5,0 \% do total, caracterizando -se como máximo local. 3.3.1.6 - Comentários sobre a correção das respostas de Formação Geral com respeito à Língua Portuguesa Na avaliação do desempenho linguístico dos participantes foram consideradas quatro competências distintas, de modo a permitir um mapeamento detalhado do domínio dos recursos disponíveis na Língua Portuguesa para a comunicação escrita formal.
Nessa situação comunicativa, as exigências para o desenvolvimento desse tipo de texto dizem respeito : à adequação da seleção vocabular, ao desenvolvimento do conteúdo, à estruturação sintática dos períodos, à organização lógica das ideias, à utilização de procedimentos de encadeamento textual e de referenciação, à obediência às exigências morfossintáticas próprias da modalidade escrita da norma-padrão, ao respeito às regras ortográficas e às regras de acentuação gráfica.
A avaliação foi realizada considerando -se os textos das respostas às duas questões de Formação Geral e os resultados obtidos em cada uma delas, como seria de se esperar, foram semelhantes.
As competências avaliadas podem ser sintetizadas nos aspectos descritos a seguir.
Aspectos textuais - estruturação textual condizente com o gênero discursivo e o modo de organização textual expositivo/argumentativo adequado ao gênero – essa competência envolve : a estruturação sintática condizente com o padrão da modalidade escrita formal da Língua Portuguesa de modo a garantir a clareza necessária ; a distribuição do conteúdo do texto em parágrafos, de modo a garantir a sua organização temática ; a utilização de operadores discursivos que contribuam para a progressão temática do texto, estabelecendo relações lógicas entre as ideias apresentadas, tanto do ponto de vista intrafrasal, como do interfrasal ; a utilização de procedimentos de referenciação lexical e pronominal que permitam a retomada de referentes textuais ; o respeito às regras de pontuação como fator de estruturação do período.
Aspectos morfossintáticos - domínio dos diferentes aspectos morfossintáticos próprios da modalidade escrita formal da norma-padrão da Língua Portuguesa – essa competência envolve : a concordância nominal, a concordância verbal, a regência nominal, a regência verbal, a flexão nominal, a flexão verbal, a correlação entre os tempos verbais, a colocação pronominal e a utilização de sinais de pontuação que contribuam para a organização lógica da frase e do texto.
Aspectos vocabulares - seleção vocabular adequada à modalidade escrita formal da Língua Portuguesa, exigida pela situação comunicativa – essa competência envolve a precisão na seleção/utilização do vocabulário relacionado à temática solicitada pela questão ; a ausência de marcas da oralidade, como termos de sentido muito genérico ( “ coisa ”, “ negócio ”, “ você ” ) e termos de registros mais informais ( como gírias, jargões, frases feitas, ditados populares, termos regionais ).
Assim, espera -se que o participante respeite a adequação vocabular não usando gírias ou expressões coloquiais, evite repetição desnecessária de palavras e utilize um vocabulário mais formal, como solicitado por um texto dissertativo.
Aspectos ortográficos - respeito às convenções ortográficas da norma-padrão da Língua Portuguesa – essa competência envolve o domínio das regras de acentuação gráfica e da grafia padrão das palavras ( com ausência de abreviaturas próprias da linguagem da internet ), de acordo com as convenções estabelecidas pela legislação em vigor e consubstanciadas no Vocabulário Ortográfico da Língua Portuguesa, editado pela Academia Brasileira de Letras. ( Vale notar que, nessa edição de 2016, não será aceita a legislação anterior, no caso das regras relativas ao uso do hífen e da acentuação gráfica ).
A escolha dessas competências para subsidiar o processo de avaliação apoia -se na concepção de que, no desempenho dos graduandos, a modalidade escrita tem apresentado uma intensa simplificação, originada no padrão da modalidade oral da Língua Portuguesa.
No caso do texto de base dissertativa, inscrito em um registro formal, a distância entre as duas modalidades é ainda maior, o que provoca situações de hipercorreção ( desvios provocados pela incorporação indevida de uma regra da norma-padrão ) e de truncamentos sintáticos ( estruturas frasais incompreensíveis devido à complexidade sintática própria da modalidade escrita ).
Para efeito de pontuação, o critério aprovado pelo Inep, e já utilizado na avaliação do desempenho linguístico nas edições do Enade desde 2013, reúne as competências acima descritas em três grupos da seguinte forma : aspectos textuais ( 40 \% ), aspectos morfossintáticos e vocabulares ( 40 \% ) e aspectos ortográficos ( 20 \% ).
A seguir, analisam -se qualitativamente os resultados da avaliação de cada um dos aspectos avaliados nos textos de resposta das questões discursivas 1 e 2 de Formação Geral.
Aspectos textuais Esta competência é a que se revelou como a mais problemática entre os participantes, porque são muitos os desvios observados.
Em relação à estrutura formal do texto produzido pelos participantes, observaram -se relações linguísticas quase agramaticais, como as estabelecidas pela sequência de gerúndios sem o apoio de um ponto de partida para a organização das informações gramaticais e semânticas, ou seja, sem oração principal.
Observou -se que uma grande parte dos participantes não organizou as ideias em parágrafos, talvez devido ao pequeno número de linhas disponibilizadas para a resposta da questão ou, quem sabe, pela suposição de que não seria necessária essa divisão por não se tratar de um texto no modelo de uma redação dissertativo-argumentativa, como solicitado nos vestibulares.
Quanto à utilização dos mecanismos de referenciação, deve -se destacar a ocorrência, em uma parte dos textos, de repetições de palavras ou expressões sem a utilização de termos sinônimos ou pronomes, como seria adequado.
Outro aspecto relevante a ser destacado é a quase total ausência de operadores argumentativos, tanto intrafrasais como interfrasais, repercutindo uma tendência atual da mídia escrita.
A respeito do uso dos sinais de pontuação, observaram -se vários desvios nos textos analisados, com a ocorrência, inclusive, de parágrafos sem marca interna de pontuação para separar os períodos.
São os seguintes os tipos de problemas encontrados, que se caracterizam como desvios ao estabelecido pela norma-padrão : a ) uso inadequado de vírgula : • ocorrência entre o sujeito e o predicado, o que é considerado inadequado por interromper a sequência natural da frase ; • ocorrência de apenas uma das vírgulas para intercalar uma palavra, uma expressão ou uma oração encaixada, quando o correto seria a demarcação com uma vírgula antes e uma depois do termo encaixado na frase ; • ocorrência da vírgula no lugar do ponto para separar ideias que constituem períodos distintos ; b ) uso inadequado de ponto e vírgula para separar elementos que deveriam ser separados por vírgula, por estarem dentro do mesmo período ; c ) ausência de ponto final para separar períodos.
Vale observar que não foi penalizada a ausência de vírgula para destacar locuções ou adjuntos adverbiais de pequena extensão deslocados de posição na frase, por ser um uso opcional.
Os problemas relativos à estruturação textual geram, portanto, os seguintes tipos de desvios : • sequência justaposta de ideias sem encaixamentos sintáticos ; • redução de estruturas subordinadas, ao lado do aumento na frequência de estruturas coordenadas e absolutas ; • redução no uso de conectores para expressar relações lógicas essenciais à construção do texto ; • emprego inadequado de operadores que não estabelecem relações lógicas coerentes entre ideias do texto ; • repetição de termos sem a utilização de procedimentos mais sofisticados de substituição ( hiperonímias, hiponímias, nominalizações, expressões metafóricas ) ; • frases fragmentadas que comprometem a estrutura lógico-gramatical ; • frases formadas apenas por oração subordinada, sem oração principal.
Aspectos morfossintáticos e vocabulares Os resultados relativos aos aspectos morfossintáticos ( concordância, regência, colocação de pronomes, flexão nominal e verbal ) evidenciam desvios recorrentes no desempenho dos participantes, relacionados ao seu processo de formação e, principalmente, a hábitos da oralidade, marcados pela informalidade.
Apesar de exemplificarem processos de mudança em curso na Língua Portuguesa, esses desvios foram penalizados por não atenderem às exigências da norma-padrão relativas ao texto escrito em situação formal, como deveria ser o texto dissertativo solicitado no Enade.
O desvio mais frequente, em relação à regência, é a falta do sinal indicativo da crase – isso revela que o usuário não tem consciência de que, sob a forma do termo “ a ”, existe a presença de uma preposição “ a ”, exigida pela regência do termo anterior, combinada a um artigo definido ( por exemplo : “ Devido as grandes guerras que estão acontecendo... ” ) Outros problemas relacionados à regência verbal e à nominal foram identificados em alguns textos dos participantes : ( a ) a ausência de preposição ou seu uso indevido antes de pronome relativo ( por exemplo : “ a região que eles moram não oferece condições de segurança ” ) ; ( b ) o uso indevido do pronome relativo ( por exemplo : “ não pensam na cidade para que eles vão ” ) ; ( c ) utilização de pronome “ onde ” deslocado de sua função locativa ).
Esses processos refletem hábitos da modalidade oral da língua, em situações de registro informal.
Apesar da possibilidade de que essa alteração de regência se generalize no padrão escrito da Língua Portuguesa, como já está ocorrendo até em textos jornalísticos, a ausência da preposição foi penalizada neste processo de avaliação.
Outro desvio observado nos textos diz respeito aos processos de concordância verbal e de concordância nominal.
Quanto à concordância de número, além das ocorrências já conhecidas de falta da marca de plural em verbos ou adjetivos, em função da ocorrência de substantivo plural na função de sujeito, observaram -se dois tipos de desvio já registrados nas pesquisas linguísticas : ( a ) a ausência de marca de plural com sujeito posposto ( por exemplo : “ Foi expulso do país os imigrantes que entraram sem autorização da polícia ” ) ; ( b ) o uso indevido da marca de plural para estabelecer concordância com o plural da locução adjetiva, apesar de o substantivo que funciona como núcleo do sintagma nominal estar no singular ( por exemplo : “ A proteção das mulheres que sofrem violência foram garantidas pela Lei Maria da Penha ” ).
Uma ocorrência generalizada foi a ausência de acento circunflexo na forma plural do presente do indicativo do verbo “ ter ”, que foi considerada como um desvio na concordância verbal e não meramente de acentuação gráfica.
Quanto à concordância de gênero, alguns casos foram observados, normalmente no âmbito de sintagmas nominais longos, em que o adjetivo está afastado do substantivo.
Quanto à questão da colocação pronominal, foram poucos os desvios observados.
Apesar de serem aspectos relacionados à oralidade, observou -se que a maioria dos participantes já incorporou regras como a não introdução da frase por um pronome oblíquo e a próclise na presença de um termo atrator.
Não se adotou, entretanto, o padrão excessivamente formal descrito pelas gramáticas normativas em relação à posição do pronome oblíquo em locuções verbais, já que esse uso está muito distante da prática cotidiana, até em textos mais formais.
Quanto aos aspectos vocabulares ( precisão semântica do vocabulário utilizado e adequação ao registro formal exigido pelo tipo de texto solicitado no enunciado ), podem ser destacados os seguintes tipos de desvios : expressões da oralidade, seleção vocabular incompatível com o contexto, gerando situações de falta de inteligibilidade ; falta de domínio de vocabulário mais abstrato e de maior complexidade, essencial ao desenvolvimento do texto de base dissertativa.
Não se observou presença de palavras de baixa calão, ao contrário do exame de 2015, em que o tema “ funk ” gerou a utilização desse tipo de vocabulário.
Várias marcas de oralidade foram identificadas, embora não em alta frequência : o uso do pronome relativo “ onde ” como relativo universal, falta de artigo definido antes de substantivo, repetição de palavras por falta de vocabulário, reduções como “ tá ”, “ pra ”, “ pro ”, “ prum ”, expressões informais.
Aspectos ortográficos : O desempenho dos participantes revelou uma diferença muito grande nos dois aspectos de caráter ortográfico analisados neste item : grafia e acentuação.
Observou -se a ocorrência de baixo índice de desvios na grafia das palavras e grande índice de desvios no uso dos sinais de acentuação gráfica.
Em vários casos, ocorre ausência completa de acentuação.
Os resultados revelam que existe uma tendência crescente entre os universitários brasileiros de eliminação da acentuação gráfica, talvez motivada pelos hábitos relacionados às redes sociais ou devido a um conhecimento ainda superficial sobre o Acordo Ortográfico da Língua Portuguesa de 1990.
Os casos mais sistemáticos de eliminação do acento indicador da sílaba tônica são : • palavras proparoxítonas ( por exemplo : “ proximos ”, “ politicos ”, “ publica ”, “ fisica ”, “ psicologica ” ) ; • palavras paroxítonas terminadas em ditongo crescente ( por exemplo : “ necessario ”, “ noticiarios ”, “ individuo ”, “ dependencia ”, “ varios ”, “ propria ”, “ violencia ”, “ ocorrencia ”, “ suicidio ” ) ; • palavras oxítonas ( por exemplo : “ ate ”, “ tambem ”, “ mes ” ) ; • palavras paroxítonas terminadas em ditongo nasal ( por exemplo : “ orgão ”, “ orfão ” ).
Quanto ao domínio das convenções relativas à grafia das palavras, pode -se concluir que os desvios estão majoritariamente relacionados a questões relativas à oralidade : • hipercorreção pela escolha de “ e ” no lugar de “ i ” ou de “ i ” no lugar de “ e ”, por influência de hábitos da oralidade : por exemplo, “ inviado ” ( no lugar de “ enviado ” ), “ entervenção ” ( no lugar de “ intervenção ” ), “ empedimento ” ( no lugar de “ impedimento ) ; • eliminação do “ r ” marcador do infinitivo verbal : por exemplo, “ esta ” ( no lugar de “ estar ” ), “ estuda ” ( no lugar de “ estudar ” ), “ migra ” ( no lugar de “ migrar ” ) ; • utilização de “ s ” no lugar de “ c ” : por exemplo, “ mensionado ” ( no lugar de “ mencionado ” ) ; “ insentivou ” ( no lugar de “ incentivou ” ) ; • utilização de “ s ” no lugar de “ x ” : por exemplo, “ espectativa ” ( no lugar de “ expectativa ”, devido à confusão com o verbo “ esperar ” ) ; e a situação contrária, “ x ” no lugar de “ s ”, “ expantoso ” ( no lugar de “ espantoso ” ) ; • utilização de “ ss ” no lugar de “ ç ” : por exemplo, “ gerassões ” ( no lugar de “ gerações ” ) ; • uso de “ am ” no lugar de “ ão ” na forma de terceira pessoa do singular do presente e do pretérito perfeito do indicativo : por exemplo, “ estam ” ( no lugar de “ estão ” ) ; • utilização de “ l ” no lugar de “ u ” em final de sílaba : por exemplo, “ tralma ” ( no lugar de “ trauma ” ) ; • outros casos de desvios de grafia relacionados à variação diastrática podem ser observados em “ estrupo ” ( no lugar de “ estupro ” ), “ subjulgadas ” ( no lugar de “ subjugadas ” ), “ precoseito ” ( no lugar de “ preconceito ” ), “ soubre ” ( no lugar de “ soube ” ), “ vecendo ” ( no lugar de “ vencendo ” ), “ indesencia ” ( no lugar de “ \textbf{indecência} ” ), “ apolojia ” ( no lugar de “ apologia ” ), “ fulga ” ( no lugar de “ fuga ” ). • inadequação no uso da maiúscula para destacar determinadas palavras-chave do texto, como “ Violência ”, “ Brasileiros ”.
Destaque -se, também, a ocorrência de participantes que grafam os textos inteiramente em caixa alta.
Vale observar que, ao contrário do esperado, não apareceram abreviaturas próprias do “ internetês ”, relacionadas ao uso de redes sociais e de e-mails.
Sua influência, portanto, parece se restringir à ausência dos acentos gráficos.
Em relação aos aspectos textuais, observa -se dificuldade na estruturação textual, o que resulta em truncamentos sintáticos e prejuízo da compreensão do texto, e desrespeito a regras essenciais de pontuação, principalmente pela dificuldade em identificar as relações lógicas que se estabelecem entre as unidades frasais.
Em relação aos aspectos morfossintáticos, domínio das regras de concordância nominal e verbal, mas desrespeito às regras de regência nominal e verbal, principalmente no que diz respeito ao uso do sinal indicativo da crase.
Em relação aos aspectos vocabulares, reduzida utilização de palavras de uso mais formal e abstrato, adequado à construção do texto de base dissertativa.
Em relação às convenções ortográficas, domínio da grafia das palavras, mas falta de atendimento à maioria das regras de acentuação gráfica. 3.3.2 - Componente de Conhecimento Específico Na parte da prova relativa às questões discursivas do Componente de Conhecimento Específico ( Tabela 3.17 ), observa -se que a Média foi mais baixa do que para as questões discursivas do Componente de Formação Geral.
Enquanto no Componente de Formação Geral a Média para estudantes de Agronomia de todo o Brasil foi 56,4, na parte de Conhecimento Específico a Média foi 51,9.
A maior Média deste componente foi obtida pelos estudantes da região Sul ( 54,6 ), e a menor, pelos da região Norte ( 48,2 ).
Quanto à variabilidade das notas, o Desvio padrão de todo o Brasil foi 18,1.
O maior Desvio padrão foi encontrado na região Norte ( 19,5 ), e o menor, na região Sul ( 16,4 ).
A maior nota Máxima, 100, foi obtida em todas as regiões, exceto na região Centro-Oeste, que obteve nota Máxima igual a 93,3.
Além disso, a nota Mínima ( 0,0 ) foi obtida por pelo menos um aluno em todas as regiões do Brasil.
A Mediana do Brasil como um todo foi 51,7.
A maior Mediana foi obtida na região Sul ( 56,7 ), e a menor, nas regiões Norte e Centro-Oeste ( 50,0 em ambas ).
O Gráfico 3.10 representa a distribuição das notas nas questões discursivas no Componente de Conhecimento Específico.
A moda desta distribuição ocorre no intervalo ( 50 ; 60 ].
A análise de cada uma destas questões será feita a seguir. 3.3.2.1 - Análise da Questão Discursiva 3 do Componente de Conhecimento Específico Na questão 3, cujos resultados aferidos encontram -se descritos na Tabela 3.18, a Média dos estudantes de todo o Brasil foi 52,9.
A menor Média nessa questão foi obtida pelos alunos da região Norte ( 45,6 ), enquanto a maior Média foi obtida na região Sul ( 58,3 ).
Quanto à variabilidade das notas, o Desvio padrão de todo o Brasil foi 28,0.
O maior Desvio padrão foi obtido na região Norte ( 29,8 ), enquanto o menor foi obtido na região Sul ( 24,6 ).
A nota Máxima, 100,0 pontos, foi alcançada por, pelo menos, um aluno de todas as regiões.
A Mediana do Brasil como um todo foi 75,0, a mesma para as regiões Nordeste, Sudeste e Sul.
A Mediana nas regiões Norte e Centro-Oeste foi 50,0.
A nota Mínima ( 0,0 ) foi obtida por, pelo menos, um estudante em todas as regiões do Brasil.
O Gráfico 3.11 mostra a distribuição das notas na questão discursiva 3, do Componente de Conhecimento Específico da área de Agronomia.
Essa distribuição é multimodal com moda principal no intervalo ( 70 ; 80 ]. 3.3.2.2 - Comentários sobre as respostas à Questão Discursiva 3 A questão 3 da parte discursiva de conhecimentos específicos da prova de Agronomia atende às exigências de formação do engenheiro agrônomo, em relação aos conteúdos, competências e habilidades presentes nas Diretrizes Curriculares da Área A formulação da questão 3 referiu -se à importância do período de carência na prevenção de prejuízos para a saúde pública, diante da ampla utilização de agrotóxicos na agricultura.
Foi também solicitado explicar o significado da expressão “ período de carência ”.
No padrão de resposta dessa questão, o estudante deveria referir -se ao período de carência como uma alternativa para amenizar o uso de agrotóxicos, assumindo que com a obediência, atendimento, e respeito ao período entre a aplicação do produto e o consumo ( colheita, comercialização, abate ou ordenha animal ), os resíduos estariam de acordo com as tolerâncias.
A questão traz um alerta para o engenheiro agrônomo sempre considerar o período de carência após o uso desses produtos.
A linguagem estava clara, direta e fácil, adequada às respostas dos estudantes.
O padrão permitia a subdivisão da pontuação e era aberto a alternativas de respostas corretas que fossem encontradas.
Como não consta, na pergunta, a relação dos agrotóxicos com os limites quantitativos de toxicidade, isso não foi abordado nas respostas fornecidas pelos estudantes.
Esclareça -se que foram aceitas como parcialmente corretas as afirmações como, por exemplo, “ o período em que o produto foi degradado ”.
O enunciado é muito claro e induz a uma resposta que envolva a saúde pública.
Considerou -se que o fundamental era o aluno demonstrar conhecimento de que é necessário um período para que o produto ( agrotóxico ) aplicado degrade ou evapore, transformando o princípio ativo em fator inerte, que não prejudique a saúde, podendo o mesmo ser consumido.
Quanto à diversidade de respostas incorretas encontradas, poderiam ser consideradas como possíveis causas : a incompreensão do enunciado ou o erro na escolha do método de solução empregado pelo estudante.
Foi constado que a linguagem adotada pelos estudantes, de um modo geral, foi adequada e de acordo com o jargão utilizado na área de agronomia, e o tempo da prova foi suficiente para a elaboração da resposta desta questão.
Cerca de metade dos estudantes que responderam à questão discursiva 3 tive um bom desempenho e a outra metade ficou dividida entre respostas fracas e medianas.
Pouco mais de 10 \% obtiveram nota zero e cerca de ¾ dos estudantes acertaram pelo menos metade da questão.
A principal abordagem utilizada pelos estudantes que obtiveram as notas mais elevadas foi relacionada ao período, ao tempo.
Os estudantes demonstraram conhecimento da necessidade de se estabelecer um período de carência para, posteriormente, consumir o produto vegetal ( poucos o relacionaram com o consumo animal ).
Os estudantes que obtiveram notas abaixo de 70,0, muitas vezes relacionavam o período de aplicação do produto na própria lavoura ou mesmo no solo, com o retorno da aplicação.
Não ocorreu uma grande diversificação de soluções para o problema proposto, visto que, na prática, não são encontradas alternativas viáveis para contornar essa situação temporal, a não ser condições especiais, como, por exemplo, lavagem de frutas e verduras em determinadas circunstâncias.
Mas isso foi pouco abordado pelos estudantes.
Foram encontradas muitas afirmações genéricas, não técnicas, como, por exemplo : que o agrotóxico faz mal, que só pode ser aplicado na área e não no produto, que é um erro aplicar agrotóxico, que só se podem usar produtos alternativos, ou que agrotóxico é um veneno e não deveria ser aplicado.
Respostas com esses tipos de conteúdo indicam que existe um déficit na formação dos agrônomos, podendo resultar na oferta de profissionais com falta de conhecimento técnico para o uso de tal produto frente aos desafios da moderna agricultura.
Pelo exposto, essa questão é extremamente relevante para o momento atual e nos faz refletir sobre a preocupação social a respeito do uso de produtos mais saudáveis para a alimentação humana e animal. 3.3.2.3 - Análise da Questão Discursiva 4 do Componente de Conhecimento Específico A Tabela 3.19 contém as informações relativas à questão 4 do conjunto de questões do Componente de Conhecimento Específico.
O desempenho dos estudantes nessa questão foi inferior ao das questões 3 e 5.
A Média geral do Brasil foi 35,6, sendo a menor Média registrada na região Norte ( 31,9 ), e a maior, na região Sul ( 37,0 ).
A nota Máxima ( 100,0 ), a Mediana ( 40,0 ) e a nota Mínima ( 0,0 ) foram atingidas por, pelo menos, um concluinte de todas as regiões.
O Gráfico 3.12 representa a distribuição de notas da questão discursiva 4, do Componente de Conhecimento Específico.
Essa distribuição é multimodal com moda principal no intervalo ( 30 ; 40 ]. 3.3.2.4 - Comentários sobre as respostas à Questão Discursiva 4 A questão 4 da parte discursiva de conhecimentos específicos da prova de Agronomia abrange técnicas que atendem às diretrizes da formação do engenheiro agrônomo, em relação aos conteúdos, competências e habilidades presentes nas Diretrizes Curriculares da Área.
A questão estava relacionada ao conteúdo de matemática – estatística – bioestatística.
O comando da questão evidencia que algum conhecimento de estatística experimental deveria ser usado, ou seja, a escolha da opção de resposta deveria estar baseada na diferença significativa expressa pelo teste Tukey 5 \%.
Entretanto, tal escolha do clone IRCA 111 também poderia se basear apenas na observação dos percentuais que se encontram na tabela do enunciado.
Isso foi observado claramente quando o estudante respondia, de forma direta, que sua escolha seria pelo “ clone IRCA 111 porque este possuía a maior ‘ média de produtividade ’, percentualmente acima dos demais clones ”.
Apesar de o estudante poder fazer a escolha correta por outro caminho, a resposta deveria exigir conhecimento de estatística experimental, o que é usual para os estudantes de agronomia familiarizados com a iniciação cientifica e a pesquisa.
O tempo da prova foi suficiente para a elaboração da resposta desta questão, bem como o espaço físico a ela destinado.
A maioria dos estudantes utilizou poucas linhas do espaço total destinado à resposta da questão.
Os estudantes que acertaram a pergunta justificaram as suas escolhas afirmando que o clone IRCA 111 diferiu significativamente dos demais clones, porque mesmo sendo igual estatisticamente, pelo teste Tukey 5 \%, ao clone PB 235, este não diferiu da testemunha.
Portanto, este foi o clone escolhido.
Outros utilizaram apenas os valores percentuais, afirmando que este clone IRCA 111 foi escolhido porque sua produtividade foi superior, de 124,82 \%, a maior em relação à testemunha.
Alguns foram sintéticos dizendo que o clone foi 24,82 \% superior à testemunha, por isso a sua escolha.
A linguagem adotada foi adequada especialmente porque se trata de uma questão objetiva e técnica.
Para esta questão, o critério de correção estabelecia três possibilidades de notas : zero, quatro ou dez, devido à objetividade da questão.
Portanto, o resultado das notas nessa questão foi bem diferente do das outras duas dissertativas.
Com esse critério, cerca de 20 \% dos estudantes tiraram nota zero, cerca de 70 \% obtiveram a nota quatro e em torno de 10 \% obtiveram a nota máxima.
Quanto à análise das respostas, é de se considerar que a abordagem dos estudantes que obtiveram uma nota fraca foi focada apenas na análise simples da produção, sem a escolha do IRCA 111.
Por exemplo : “ o clone PB 235 foi escolhido porque ele teve a maior produtividade acima da testemunha ” ; “ os clones IRCA 111 e PB 235 foram superiores aos demais ” ; “ o clone PB 28/29, por ser o clone que chegou mais perto da testemunha ” ; “ a seringueira recomendada é a PB 235, pois teve maior produtividade que as demais ”.
Os estudantes que obtiveram notas médias definiram o clone IRCA 111 usando as seguintes abordagens : “ o IRCA apresentou a maior média de produtividade ” ; “ o clone IRCA 111 apresentou maior adaptação em relação aos outros clones e testemunha ” ; “ o mais recomendado é o IRCA 111, pois o mesmo apresentou a maior produtividade se comparado à testemunha ”.
Os estudantes com as notas consideradas boas indicaram o IRCA 111 com base em análise estatística usando a seguinte abordagem : “ de acordo como os resultados estatísticos apresentados na tabela, o clone IRCA 111 é o que melhor apresentou respostas satisfatórias e diferenças significativas em relação à testemunha ” ; “ é recomendado para o plantio o clone IRCA 111, que difere estatisticamente da testemunha que apresenta a maior média da produção/ha/ano bem superior à testemunha ”.
O principal erro cometido foi não utilizarem a análise estatística como ferramenta de separação das médias de produção.
A questão pode ser considerada a mais direta entre as discursivas da prova da Agronomia.
Ela aferiu competência e habilidades dos estudantes de Agronomia em conteúdo fundamental para a separação dos graus entre os níveis médio e superior.
Um engenheiro agrônomo deve utilizar as ferramentas que sua profissão oferece para melhorar a produtividade ou produção.
Daí porque o conhecimento de matemática juntamente com a discussão de dados utilizando a bioestatística ou a estatística torna -se essencial para a resolução da questão. 3.3.2.5 - Análise da Questão Discursiva 5 do Componente de Conhecimento Específico A Tabela 3.20 contém as informações relativas à questão 5 do conjunto do Componente de Conhecimento Específico.
O desempenho dos estudantes de todo o Brasil nesta questão foi mais alto dentre as três questões discursivas desse componente.
A nota Média dos estudantes de todo o Brasil foi 67,3.
A maior Média foi registrada na região Sul ( 68,4 ), enquanto a menor Média foi registrada na região Centro-Oeste ( 64,6 ).
Quanto à variabilidade das notas, o Desvio padrão dos alunos do Brasil, como um todo, foi 25,9.
Enquanto o maior desvio foi encontrado na região Norte ( 28,2 ), o menor foi encontrado na região Sul ( 24,1 ).
A Mediana para o Brasil foi 80,0, a mesma obtida em todas as regiões, exceto a região Centro-Oeste, que obteve Mediana igual a 60,0.
Para o conjunto de alunos de Agronomia do Brasil, a nota Máxima foi 100,0, e a nota Mínima ( 0,0 ).
Os mesmos limites foram encontrados em todas as regiões.
O Gráfico 3.13 representa a distribuição das notas da questão discursiva 5 do Componente de Conhecimento Específico.
Essa distribuição também é multimodal com moda principal no intervalo ( 70 ; 80 ], seguida pelo intervalo ( 50 ; 60 ].
Quanto à variabilidade das notas, o Desvio padrão dos alunos do Brasil, como um todo, foi 25,9.
Enquanto o maior desvio foi encontrado na região Norte ( 28,2 ), o menor foi encontrado na região Sul ( 24,1 ).
A Mediana para o Brasil foi 80,0, a mesma obtida em todas as regiões, exceto a região Centro-Oeste, que obteve Mediana igual a 60,0.
Para o conjunto de alunos de Agronomia do Brasil, a nota Máxima foi 100,0, e a nota Mínima ( 0,0 ).
Os mesmos limites foram encontrados em todas as regiões.
O Gráfico 3.13 representa a distribuição das notas da questão discursiva 5 do Componente de Conhecimento Específico.
Essa distribuição também é multimodal com moda principal no intervalo ( 70 ; 80 ], seguida pelo intervalo ( 50 ; 60 ]. 3.3.2.6 - Comentários sobre as respostas à Questão Discursiva 5 A canção “ Xote Ecológico ” sugere que a intervenção do homem na natureza está destruindo o meio ambiente.
O enunciado estava claro, contextualizou a intervenção do homem na natureza, na agricultura, solicitando ao estudante uma visão crítica que, ao mesmo tempo, deveria apresentar soluções.
A linguagem de arte poética utilizada facilitou a compreensão da grave situação do meio ambiente no mundo, nos dias atuais.
Os itens ‘ a ’ e ‘ b ’ da questão eram objetivos e abrangentes, o que possibilitou inúmeras respostas válidas.
O conteúdo está dentro das Diretrizes do programa de Agronomia, envolvendo situação do contexto agrícola.
O espaço disponível para a resolução da questão foi adequado, proporcional ao tempo de prova, mesmo considerando a abrangência das perguntas.
Tal abrangência possibilitou uma grande variedade de respostas, a maioria com foco correto.
Esta questão foi a de menor dificuldade entre as três questões discursivas da prova de Agronomia, pois tratou de tema atual, presente no cotidiano da mídia e constante nos programas de graduação em agronomia.
Em relação ao padrão de resposta e às respostas fornecidas pelos estudantes, é de se notar que a maioria dos estudantes respondeu que a produção de alimentos, de forma convencional, nos modos de hoje, sem preocupação com a conservação ambiental, conduziria ao esgotamento produtivo do ambiente.
Foram citados exemplos do excesso de desmatamento, de uso de agrotóxico e de adubos, de monocultura, de queimadas, etc.
Muitas respostas citaram alternativas de manejo que foram contempladas no padrão de resposta, como plantio direto, rotação de culturas, pastoreio rotativo, cultivos agroecológicos, controle biológico, manejo racional da água, descarte correto de embalagens, entre outros.
As respostas fracas usaram termos como “ primeiro deve ser pensada a poluição na cidade... ” ; “ falta incentivo do governo e dos professores... ” “ não se deve usar agrotóxicos nunca... ”.
Os estudantes, que responderam de forma correta às perguntas, forneceram, de alguma forma, respostas que estavam de acordo com o esperado pelo padrão de resposta da questão.
A maioria das respostas foi sintética, precisa, dizendo ser necessário maior interação entre lavoura e pecuária, fazer uso de plantio direto, preservar matas ciliares, utilizar praticas ecológicas.
A diversidade de resposta ocorreu de forma positiva, sempre alertando para a preocupação ambiental e, portanto, o percentual de notas baixas foi muito pequeno.
A linguagem utilizada pelos estudantes foi adequada, recorrendo a termos específicos da área de agronomia.
Quanto ao desempenho do total de estudantes que responderam à questão, menos de 1 \% recebeu nota zero e acima de 80 \% receberam notas acima de 6,0, o que caracteriza um bom desempenho na questão.
Em relação às respostas fornecidas pelos estudantes, foram encontradas respostas bem elaboradas, com percepção correta da realidade do homem na terra, hoje em dia, embora, às vezes, apresentando soluções simplistas, sem indicação de alternativas que deveriam ser tomadas, como na seguinte resposta : “ Hoje temos infinitas ações para diminuir o impacto ambiental e melhorar a qualidade de vida, há anos atrás tínhamos uma cultura de desmatar para plantar, hoje essa técnica foi abandonada por prejuízos ao meio ambiente, causando seca, degradação ambiental, entre outros.
Com avanço das tecnologias e estudos obtivemos técnicas que hoje conseguimos o aumento da produção visando o manejo e técnicas de lavoura-pecuária-floresta... ”.
Esta resposta é um bom exemplo de um estudante que apenas alertou para o problema, mas não relatou quais técnicas para diminuí -lo, não respondeu parte importante da pergunta, de modo a demonstrar conhecimento específico de um engenheiro agrônomo.
Outros preferiram simplificar ao máximo a resposta dizendo que a música só diz respeito à intervenção do homem, também não respondendo à questão.
A principal abordagem adotada para responder aos itens ‘ a ’ e ‘ b ’ foi a questão ambiental, mostrando que os estudantes estão preocupados com a realidade.
O alto índice de acertos da questão reflete isso.
A questão não exigiu que fosse explicitada a metodologia para realizar o plantio direto, por exemplo, bastava citá -la.
Não foram, também, utilizados exemplos práticos nas respostas, referenciando a resposta com o que já aconteceu no Brasil, por exemplo, com a mata atlântica, com o cerrado, ou na Amazônia.
Apesar da diversidade de respostas, a maioria entendeu a mensagem da música, e até usou citações de sua letra nas suas respostas.
Observando as respostas, ficou claro que a comunidade envolvida na área de agronomia, professores, pesquisadores e estudantes necessitam investir mais em ciências neste campo do conhecimento agronômico, buscando novos conhecimentos da interação produção de alimentos x ambiente.
Na prática, existem poucas alternativas para um manejo ecológico e com produtividade econômica na agricultura brasileira, principalmente para o caso das grandes culturas, anuais.
Os manejos como o MIP e o uso do controle biológico têm, ao longo dos últimos anos, sido abandonados, muito em função do desconhecimento, o que pode ser fruto da falta do ensinamento e pesquisa nos centros de ensino.
A diminuição dos preços de alguns produtos agrotóxicos e a facilidade de encontrá -los no mercado, muito diferente das décadas de 70 e 80, por exemplo, deve ter facilitado o seu uso indiscriminado.
Por essa razão, faz -se necessário mais conhecimento envolvendo o aprendizado de mais alternativas de manejo, oferecendo opções ao uso destes produtos.
A falta de conhecimento explica o fato de os estudantes responderem de forma genérica, recorrendo a saberes do senso comum. 3.3.3 - Considerações Finais As questões discursivas 3, 4 e 5 da prova abordaram assuntos de importância atual na agricultura e avaliaram o grau de conhecimento dos estudantes frente a temas de amplitude nacional e de importância na vida profissional do Engenheiro Agrônomo.

% matched lemmas: aportar, autoestima, homicídio, humilhação, indecência
\end{textsample}
